\chapter{Introduction}

\begin{summarybox}{Synthèse}
    Démonstration des nouveaux réglages : numérotation continue des figures, unités en km/h, et profondeur des titres.
\end{summarybox}

\section{Niveaux de Titres (Section)}
Ceci est le texte standard sous une section.

\subsection{Titre de niveau 2 (Subsection)}
On descend dans la hiérarchie.

\subsubsection{Titre de niveau 3 (Subsubsection)}
Encore un niveau en dessous.

\paragraph{Titre de niveau 4 (Paragraph)} 
Notez que le texte commence juste après le titre sur la même ligne (style "runin"). C'est très utile pour les définitions ou les points techniques.

\subparagraph{Titre de niveau 5 (Subparagraph)} 
Le niveau le plus bas, souvent en italique ou gris pour des nuances fines.

\section{Démonstration Siunitx (Mise à jour)}

Voici la comparaison des affichages avec l'option \texttt{per-mode = symbol} activée :

\begin{itemize}
    \item Vitesse : \qty{90}{\kilo\meter\per\hour} (au lieu de km.h\textsuperscript{-1})
    \item Accélération : \qty{9.81}{\meter\per\second\squared} (affiche m/s\textsuperscript{2})
    \item Débit : \qty{50}{\liter\per\minute}
\end{itemize}

\section{Test Numérotation Figure}

La figure ci-dessous devrait être numérotée "Figure 1" et non "Figure 1.1".

\begin{figure}[hbt!]
    \centering
    \includegraphics[height=8cm]{img/logos/Logo_LICIT-ECO7_vertical.png}
    \caption{Exemple de figure avec numérotation continue}
    \label{fig:test_num}
\end{figure}

Comme on le voit sur la \cref{fig:test_num}, le numéro est global.